problem:  
Let \( M^n \subset \mathbb{H}^{n+1} \) be a complete, properly embedded minimal hypersurface that is *asymptotically regular* in the following sense: in the conformal compactification of \(\mathbb{H}^{n+1}\) by a defining function \(\rho\), the induced metric on \(M\) admits a Fefferman–Graham-type expansion  
\[
g_M = \rho^{-2}(h_0 + \rho^2 h_2 + \cdots + \rho^n h_n + o(\rho^n)),
\]
where \(h_0\) is the induced metric on the asymptotic boundary \(\partial_\infty M\subset S^n_\infty\).  
Denote the renormalized volume (finite part of the regularized volume expansion) by \(V_{\mathrm{ren}}(M)\).

The Graham–Witten theory implies that \(V_{\mathrm{ren}}(M)\) is a conformal invariant of the **embedding pair** \((\partial_\infty M,[h_0])\), generally involving local conformal invariants of the embedded submanifold \(\partial_\infty M \subset S^n_\infty\).

**Problem (Conjecture):**  
Suppose \(M^n\) is asymptotically umbilic in the sense that the trace-free second fundamental form \(A^\circ\) satisfies \( |A^\circ|(x) = O(e^{-\alpha\, d_{\mathbb{H}^{n+1}}(x,x_0)}) \) for some \(\alpha>0\).  
Does this extra geometric decay imply that the *renormalized volume* \(V_{\mathrm{ren}}(M)\) becomes a *purely intrinsic* conformal invariant of the boundary metric \([h_0]\); that is,  
\[
V_{\mathrm{ren}}(M) = \mathcal{F}([h_0]),
\]
independent of the embedding of \(\partial_\infty M\) in \(S^n_\infty\)?  
Equivalently, does exponential asymptotic umbilicity “erase” all extrinsic contributions to the holographic coefficients?

Why is it a "good" problem:  
This refined conjecture isolates a subtle boundary rigidity phenomenon: whether rapid extrinsic decay transforms a mixed (intrinsic + extrinsic) conformal invariant into a purely intrinsic one. It stays within established analytic frameworks (Fefferman–Graham expansions, Graham–Witten renormalized volume) while probing their limits.

- It is **mathematically coherent and precise**: the hypotheses are standard in conformally compact geometry, and the expansion used is well‑defined.  
- It addresses a **deep structural question** about how bulk decay of extrinsic curvature controls the asymptotic invariants at infinity.  
- It **bridges distinct areas**—minimal surface theory, conformal geometry, and holographic renormalization—seeking to determine when the boundary conformal class alone suffices to encode bulk renormalized quantities.  
- It is **open**: current formulae for \(V_{\mathrm{ren}}\) involve extrinsic curvature terms, and it is not known whether these vanish or simplify under exponential umbilicity decay. Resolving this would require new harmonic‑analytic estimates and boundary regularity results.  

Thus the problem is conceptually clean, technically challenging, and connects analytic, geometric, and conformal perspectives in an unexplored way—qualifying as a genuinely “good” research‑level mathematical problem.